\documentclass[12pt,letterpaper]{article}



%%
%% Required packages:
%%
\usepackage{etex}
\usepackage{amsmath}
\usepackage{amssymb}
\usepackage{amstext}
\usepackage{amsthm}
\usepackage{fancyhdr,fancybox}
\usepackage{mathrsfs} % need for \mathscr command
\usepackage{multicol}
\usepackage[normalem]{ulem}
%\usepackage{pictex,pictexplus}
% \usepackage{pictexwd}
\usepackage{rotating}
\usepackage{url}
\usepackage{xfrac}
\usepackage{booktabs}
\usepackage{color,soul}
\usepackage{comment}

\usepackage{hyperref}
\hypersetup{
    colorlinks=true,     % false: boxed links; true: colored links
    linkcolor=blue,      % color of internal links
    citecolor=blue,      % color of links to bibliography
    filecolor=blue,      % color of file links
    urlcolor=blue        % color of external links
}


\usepackage{tikz}
\usetikzlibrary{backgrounds,calc}

%%
%% Common formatting commands
%%
\setlength{\parindent}{0in}
\setlength{\textwidth}{7in}
\setlength{\evensidemargin}{-0.25in}
\setlength{\oddsidemargin}{-0.25in}
\setlength{\parskip}{.5\baselineskip}
\setlength{\topmargin}{-0.5in}
\setlength{\textheight}{9in}

%               Problem and Part
\newcounter{problemnumber}
\newcounter{partnumber}
\newcounter{subpartnumber}
\newcommand{\Problem}{%
  \refstepcounter{problemnumber}%
  \setcounter{partnumber}{0}%
  \setcounter{subpartnumber}{0}%
  \item[\makebox{\hfil\textbf{\theproblemnumber.}\hfil}]%
  }
\newcommand{\InProblem}[2][1.6]{%
  \refstepcounter{problemnumber}%
  \setcounter{partnumber}{0}%
  \setcounter{subpartnumber}{0}%
  \textbf{\theproblemnumber.}\ \ \parbox[t]{#1in}{#2}%
}
\newcommand{\InSmallProblem}[1]{\InProblem[1.05]{#1}}
\newcommand{\NoProblem}[1]{%
  \mbox{\hfil\textbf{#1}\hfil}%
  }
\newcommand{\UnProblem}{%
  \refstepcounter{problemnumber}%
  \setcounter{partnumber}{0}%
  \setcounter{subpartnumber}{0}%
  \mbox{\hfil\textbf{\theproblemnumber.}\hfil}%
  }
\newcommand{\InFlexProblem}[1]{%
  \refstepcounter{problemnumber}%
  \setcounter{partnumber}{0}%
  \setcounter{subpartnumber}{0}%
  \textbf{\theproblemnumber.}\ \ #1%
}
%%  Part:
\newcommand{\Part}{%
  \renewcommand{\thepartnumber}{(\alph{partnumber})}%
  \refstepcounter{partnumber}
  \setcounter{subpartnumber}{0}%
  \item[(\alph{partnumber})]%
}
\newcommand{\InPart}[2][1.75]{%
  \renewcommand{\thepartnumber}{(\alph{partnumber})}%
  \refstepcounter{partnumber}%
  \setcounter{subpartnumber}{0}%
  (\alph{partnumber})\ \ \parbox[t]{#1in}{#2}%
}
\newcommand{\UnPart}{%
  \refstepcounter{partnumber}%
  \renewcommand{\thepartnumber}{(\alph{partnumber})}%
  \setcounter{subpartnumber}{0}%
  (\alph{partnumber})%
}
\newcommand{\InLargePart}[1]{\InPart[3.05]{#1}}
\newcommand{\InFullPart}[1]{\InPart[6.2]{#1}}
\newcommand{\InSmallPart}[1]{\InPart[1.05]{#1}}
%% SubPart:
\newcommand{\SubPart}{%
  \refstepcounter{subpartnumber}%
  \item[(\roman{subpartnumber})]%
  }
\newcommand{\InSubPart}[2][2.25]{%
  \stepcounter{subpartnumber}%
  (\roman{subpartnumber})\ \ \parbox[t]{#1in}{#2}%
}
\newcommand{\InSmallSubPart}[1]{\InSubPart[1.5]{#1}}
\newcommand{\InTinySubPart}[1]{%
  \stepcounter{subpartnumber}%
  (\roman{subpartnumber})\ \ \makebox{#1}%
}
\newcommand{\InMySubPart}[2][2.25]{%
  \stepcounter{subpartnumber}%
  \parbox[t]{#1in}{\makebox[0.3in][l]{(\roman{subpartnumber})}\ \ {#2}}%
}
\newcommand{\TrueFalse}{\textbf{T}\ \ \textbf{F}\ \ }
\newcommand{\TFPart}[1]{% 
  \stepcounter{partnumber}%
  \setcounter{subpartnumber}{0}%
  \item[(\alph{partnumber})] \TrueFalse\parbox[t]{5.5in}{#1}}

\newcommand{\Points}[1]{(#1 points) \ }
\newcommand{\Point}[1]{(#1 point) \ }


%% For Answers:
\newcounter{answernumber}
\newcommand{\Answer}{%
  \refstepcounter{answernumber}%
  \setcounter{partnumber}{0}%
  \setcounter{subpartnumber}{0}%
  \item[\makebox{\hfil\textbf{\theanswernumber.}\hfil}]%
  }


% Setting up the headers for the first page
%\chead{} % will default to what is typed in the file.
\fancyhead[L]{\textsc{Math 22a}}
\fancyhead[R]{\textsc{Fall 2024}}
\fancyfoot[R]{
	\vspace*{\fill}\footnotesize\textcopyright{} \emph{Harvard Math 22a}	
}
\renewcommand{\headrulewidth}{0pt}

%\fancyhead[C]{}
%	\chead{}	
%	\lhead{}
%	\rhead{}
%	\cfoot{}


% Putting copyright on the footer of each page
%\fancypagestyle{secondpage}{
%	\fancyhead[C]{TESTing again} % change this to be blank
%		%\chead{}	
%	\fancyhead[L]{bears} % change this to be blank
%	\fancyhead[R]{dogs} % change this to be blank
%	\fancyfoot[R]{ %keeps the footer the same
%		\vspace*{\fill}\footnotesize\textcopyright{} \emph{Harvard Math 22a}	
%	}
%}
%\renewcommand{\headrulewidth}{0pt}
%\pagestyle{fancy}
\pagestyle{empty}


%\lhead{}
%\rhead{}
%\rfoot{\vspace*{\fill}\footnotesize\textcopyright{} \emph{Harvard Math 22a}}
%\renewcommand{\headrulewidth}{0pt}
%\pagestyle{fancy}




%%
%% Common shortcut commands:
%%
\newcommand{\ds}{\displaystyle}
\newcommand{\sgn}{\operatorname{sgn}}
\renewcommand{\Pr}{\operatorname{P}}
\newcommand{\CPr}[3][{}]{\Pr_{\text{#1}}\left(#2\,|\,#3\right)}

\newcommand{\C}{\mathbb{C}}
\newcommand{\R}{\mathbb{R}}
\newcommand{\Z}{\mathbb{Z}}
\newcommand{\N}{\mathbb{N}}
\newcommand{\Q}{\mathbb{Q}}
\newcommand{\D}{\mathbb{D}}

\newcommand{\rref}{\operatorname{rref}}
\newcommand{\im}{\operatorname{im}}
\newcommand{\rank}{\operatorname{rank}}
\newcommand{\diag}{\operatorname{diag}}
\newcommand{\Span}{\operatorname{span}}
%\renewcommand{\vec}[1]{\mathbf{#1}}
\newcommand{\charpoly}{\mathscr{P}}
\newcommand{\nicefrac}[2]{\sfrac{#1}{#2}} % using xfrac package
\newcommand{\topology}{\mathcal{T}}
\newcommand{\Inn}{\operatorname{Inn}}

\newcommand{\blankbox}[2]{\fbox{\rule{#1}{0in}\rule{0in}{#2}}}

\newenvironment{myindentpar}[1]%
 {\begin{list}{}%
         {\setlength{\leftmargin}{#1}
          \setlength{\rightmargin}{\leftmargin}}%
         \item[]%
 }
 {\end{list}}
\newtheorem{theorem}{Theorem}
\newtheorem*{NStheorem}{Theorem}
\newtheorem{cor}[theorem]{Corollary}
\newtheorem*{claim}{Claim}
\newtheorem{defn}{Definition}

\newenvironment{amatrix}[1]{%
  \left(\begin{array}{@{}*{#1}{c}|c@{}}
}{%
  \end{array}\right)
}

%--------grstep
% For denoting a Gauss' reduction step.
% Use as: \grstep{\rho_1+\rho_3} or \grstep[2\rho_5 \\ 3\rho_6]{\rho_1+\rho_3}
\newcommand{\grstep}[2][\relax]{%
   \ensuremath{\mathrel{
       {\mathop{\longrightarrow}\limits^{#2\mathstrut}_{
                                     \begin{subarray}{l} #1 \end{subarray}}}}}}
\newcommand{\swap}{\leftrightarrow}


\usetikzlibrary{arrows}
\usepackage{wrapfig}

\makeatletter
\renewcommand*\env@matrix[1][*\c@MaxMatrixCols c]{%
	\hskip -\arraycolsep
	\let\@ifnextchar\new@ifnextchar
	\array{#1}}
\makeatother

\def\env@matrix{\hskip -\arraycolsep
	\let\@ifnextchar\new@ifnextchar
	\array{*\c@MaxMatrixCols c}}

\chead{\Large\textbf{Eigenvalues and Eigenvectors, $\vec\S$5.1, 5.4}}% 

\begin{document}
\thispagestyle{fancy}

Recall from last class:

\begin{itemize}
\item {\bf Theorem:} Let $\mathcal{B}=\{b_1, \dots, b_n\}$ and $\mathcal{C}=\{c_1, \dots, c_n\}$ be bases of a vector space $V$. Then there exists a unique $n\times n$ matrix $P$ such that $[x]_{\mathcal{C}}=P [x]_{\mathcal{B}}$.
\item We call $P$ the {\bf change of coordinates matrix from $\mathcal{B}$ to $\mathcal{C}$}. In general $$P= [[b_1]_\mathcal{C} \; \dots \; [b_n]_\mathcal{C}].$$ Also $P^{-1}$ is the change of coordinates matrix from $\mathcal{C}$ to $\mathcal{B}$.
\item Let $V$ be a vector space with basis $\mathcal{B}=\{b_1, \dots, b_n\}$ and $W$ be a vector space with basis $\mathcal{C}=\{c_1, \dots, c_m\}$. Given a linear transformation $T: V\to W$, we naturally seek an associated matrix to represent $T$ using the bases; namely we want an $m\times n$ matrix $M$ such that $[T(x)]_\mathcal{C}=M [x]_\mathcal{B}$ for all $x \in V$. We call $M$ the {\bf matrix of $T$ relative to bases $\mathcal{B}$ and $\mathcal{C}$}. In general $M= [[T(b_1)]_\mathcal{C} \; \dots \; [T(b_n)]_\mathcal{C}]$.
\item {\bf Diagonal Representation Theorem:} Suppose $A= P D P^{-1}$ where $D$ is a diagonal matrix and $T(x)=Ax$. If $\mathcal{B}$ is the basis for $\mathbb{R}^n$ formed by the columns of $P$, then $D$ is the matrix of $T$ relative to $\mathcal{B}$. 
\item If $A$ and $B$ are $n\times n$ matrices, then $A$ is {\bf similar} to $B$ if there exists an invertible matrix $P$ such that $A=PB P^{-1}$.
\end{itemize}

\newpage

\begin{defn}
If $A$ is an $n \times n$ matrix, then $\lambda \in \mathbb{R}$ is an {\bf eigenvalue} of $A$ if there exists a non-zero vector $v$, called an {\bf eigenvector} corresponding to $\lambda$, such that $$Av=\lambda v.$$
\end{defn}

\begin{enumerate}

\Problem Let $A=\begin{bmatrix} 1 & 3\\ 4 & 2\\ \end{bmatrix}$ and $v=\begin{bmatrix} 1 \\ -1\\ \end{bmatrix}$. Determine if $v$ is an eigenvector of $A$ and find the corresponding eigenvalue.

\vfill

\newpage

\begin{defn} If $\lambda$ is an eigenvalue of an $n \times n$ matrix $A$, then the set of all solutions $v$ to $Av = \lambda v$ is the {\bf eigenspace} of $A$ corresponding to $\lambda$.
\end{defn}

\Problem Let $A=\begin{bmatrix} 4 & -1 & 6\\ 2 & 1 & 6\\ 2 & -1 & 8\\ \end{bmatrix}$, then $\lambda=2$ is an eigenvalue. Find the corresponding eigenspace.

\vfill

{\bf Theorem 1:} The eigenvalues of a triangular matrix are the diagonal entries.

\Problem Prove the theorem.

\vfill

\newpage

{\bf Theorem 2:} If $v_1, \dots, v_p$ are eigenvectors corresponding to distinct eigenvalues $\lambda_1, \dots, \lambda_p$ of an $n\times n$ matrix $A$, then $v_1, \dots, v_p$ are linearly independent.

\Problem Prove the theorem

\vfill

\Problem For a general $n \times n$ matrix $A$, how can we find the eigenvalues?

\vfill

\Problem Let $A=\begin{bmatrix} 3 & 2\\ 3 & 8\\ \end{bmatrix}$ Find the eigenvalues of $A$ and corresponding eigenspaces.

\vfill

\newpage

\Problem Prove that $\lambda$ is an eigenvalue of $A$ if and only if $\lambda$ is an eigenvalue for $A^{T}$.

\vfill

\begin{defn} If $A$ is an $n\times n$ matrix, then the characteristic polynomial of $A$, denoted by $p_A(\lambda)$, is given by $$p_A(\lambda)=\text{det}(A-\lambda I_n).$$
\end{defn}

{\bf Note:} $\lambda$ is an eigenvalue of $A$ if and only if $\lambda$ is a zero of the characteristic polynomial.

\Problem Let $A=\begin{bmatrix} 1 & 2 & 3 & 4\\ 
0 & 5 & 6 & 6\\
0 & 0 & 5 & 5\\
0 & 0 & 0 & 2\\
\end{bmatrix}$ Find the characteristic polynomial of $A$.

\vfill 

\newpage

\Problem If $A$ and $B$ are similar matrices, then they have the same characteristic polynomials and hence the same eigenvalues (with the same algebraic multiplicities). 

\vfill

\Problem {\bf True or False.}

\Part If two matrices have the same eigenvalues, then they are similar.

\vfill

\Part If two matrices are row equivalent, then they have the same eigenvalues.

\vfill

\Part If two matrices are row equivalent, then they are similar.

\vfill

\end{enumerate}

\begin{comment}
\newpage
\chead{\Large\textbf{Eigenvalues and Eigenvectors -- Answers / Solutions}}%
\lhead{}
\rhead{}
\thispagestyle{fancy}

\begin{enumerate}
  \Answer  Let $A=\begin{bmatrix} 1 & 3\\ 4 & 2\\ \end{bmatrix}$ and $v=\begin{bmatrix} 1 \\ -1\\ \end{bmatrix}$, then $$Av = \begin{bmatrix} 1 & 3\\ 4 & 2\\ \end{bmatrix}\begin{bmatrix} 1 \\ -1\\ \end{bmatrix} = \begin{bmatrix} -2\\ 2\\ \end{bmatrix} = (-2) v.$$ Thus multiplying $v$ by $A$ is just scaling $v$ by $-2$, so $v$ is an eigenvector corresponding to the eigenvalue of $-2$.
    \Answer Recall the eigenspace corresponding to the eigenvalue of $\lambda=2$ is given by $E_2=N(A-2 I_3)$. Note $$A-2I_3 =\begin{bmatrix} 2 & -1 & 6\\ 2 & -1 & 6\\ 2 & -1 6\\ \end{bmatrix} \sim \begin{bmatrix} 2 & -1 & 6\\ 0 & 0 & 0\\ 0 & 0 & 0\\ \end{bmatrix}.$$ Thus $$N(A-2I_3) =\text{Span} \left\{\begin{bmatrix} 1\\ 2\\ 0\\ \end{bmatrix}, \begin{bmatrix} -3 \\ 0\\ 1\\ \end{bmatrix} \right\}.$$
     \Answer 
     
     \begin{proof} Let $A$ be an upper-triangular matrix, then we want to know what $\lambda \in \mathbb{R}$ allow for non-trivial solutions to $(A-\lambda I_n)x=0$. Thus we have
     $$A-\lambda I_n =\begin{bmatrix} a_{11}-\lambda & & *\\
       & \ddots & \\
      0 & \cdots & a_{nn}-\lambda \\ \end{bmatrix}.$$
      
      Thus $(A-\lambda I_n)x = 0$ has non-trivial solutions if and only if $A-\lambda I_n$ is not invertible by the Invertible Matrix Theorem. Again by the Invertible Matrix Theorem, $A-\lambda I_n$ is not invertible is equivalent to $\text{det}(A-\lambda I_n)=0$. We proved the determinant of an upper-triangular matrix is the product of the diagonal entries. Thus we have $$0=\text{det}(A-\lambda I_n)=(a_{11}-\lambda)\dots (a_{nn}-\lambda).$$ Thus $\lambda$ is an eigenvalue for $A$ if and only if $\lambda$ is one of $a_{11}, \dots, a_{nn}$.
     \end{proof}
      \Answer 
      
      \begin{proof} We prove the result by contradiction. Suppose $v_1, \dots, v_p$ are linearly dependent. Since $v_1\neq 0$, there is a smallest $r>1$ such that $$v_r = c_1 v_1+ \dots + c_{r-1} v_{r-1}.$$
Applying $A$ to both sides and simplifying, we get the following
\begin{align*}
Av_r & = A( c_1 v_1+ \dots + c_{r-1} v_{r-1})\\
A v_r & = c_1 Av_1+ \dots + c_{r-1} Av_{r-1}\\
\lambda_r v_r &= c_1 \lambda_1 v_1+ \dots + c_{r-1} \lambda_{r-1} v_{r-1}\\
\lambda_r (c_1 v_1+ \dots + c_{r-1} v_{r-1}) &= c_1 \lambda_1 v_1+ \dots + c_{r-1} \lambda_{r-1} v_{r-1}.
\end{align*}
Rearranging the equation, we get $$0=(c_1\lambda_1 -\lambda_r c_1) v_1 +\dots +(c_{r-1} \lambda_{r-1}-\lambda_r c_{r-1})v_{r-1}.$$

Since $v_1, \dots, v_{r-1}$ are linearly independent, we get $$(c_1\lambda_1 -\lambda_r c_1) = \dots =(c_{r-1} \lambda_{r-1}-\lambda_r c_{r-1})=0.$$ Since $\lambda_1,\dots, \lambda_r$ are distinct, we must have $c_1=\dots=c_{r-1}=0$. This implies $v_{r}=0$, but eigenvectors must be non-zero. Thus we have a contradiction.
      \end{proof}
      \Answer We've seen several times now that $\lambda$ is an eigenvalue for $A$ if and only if $A-\lambda I_n$ is not invertible. Thus we can find the eigenvalues of $A$ by finding all values of $\lambda$ such that $\text{det}(A-\lambda I_n) =0$.
      
       \Answer Let $A=\begin{bmatrix} 3 & 2\\ 3 & 8\\ \end{bmatrix}$, then $$p_A(\lambda) =\det \begin{bmatrix} 3-\lambda & 2\\ 3 & 8-\lambda\\ \end{bmatrix}=(3-\lambda)(8-\lambda)-6=(\lambda -9)(\lambda -2).$$ Thus the eigenvalues of $A$ are $2$ and $9$.
       
\underline{$\lambda =2$} Here we have $A-2I_2 =\begin{bmatrix} 1 & 2\\ 3 & 6\\ \end{bmatrix}\sim \begin{bmatrix} 1 & 2\\ 0 & 0 \\ \end{bmatrix}$. Thus $$E_2 =N(A-2I_2) =\text{Span}\left\{\begin{bmatrix}-2\\ 1\\ \end{bmatrix} \right\}.$$

\underline{$\lambda =9$}  Here we have $A-9I_2 =\begin{bmatrix} -6 & 2\\ 3 & -1\\ \end{bmatrix}\sim \begin{bmatrix} 3 & -1\\ 0 & 0 \\ \end{bmatrix}$. Thus $$E_9 =N(A-9I_2) =\text{Span}\left\{\begin{bmatrix}1\\ 3\\ \end{bmatrix} \right\}.$$

Thus $$A=\begin{bmatrix} 1 & -2\\ 3 & 1\\ \end{bmatrix} \begin{bmatrix} 9 & 0 \\ 0 & 2 \\ \end{bmatrix} \begin{bmatrix} 1 & -2\\ 3 & 1\\ \end{bmatrix}^{-1}.$$
        \Answer 
        
        \begin{proof}
        Since $p_{A^{T}} =\det(A^{T}-\lambda I_n) =\det (A-\lambda I_n)^{T} =\det (A-\lambda I_n) =p_{A}$, we get $A$ and $A^{T}$ have the same characteristic polynomials. Thus they have the same eigenvalues with the same algebraic multiplicities.
        \end{proof}
        
         \Answer From the previous theorem, the characteristic polynomial is given by $$p_{A}(\lambda)= (\lambda-1)(\lambda-5)^2(\lambda-2).$$ Thus 1 and 2 are eigenvalues with algebraic multiplicities 1 and 5 is an eigenvalue with algebraic multiplicity 2. 
         
           \Answer 
           
           \begin{proof}
           Since $A$ and $B$ are similar, there exists an invertible matrix $P$ such that $A=PBP^{-1}$. Then 
           $$p_{A}(\lambda)= \det(A-\lambda I_n)=  \det(PBP^{-1}-\lambda PP^{-1})=\det{P} p_{B}(\lambda) \det{P^{-1}} = p_{B}(\lambda).$$ Thus $A$ and $B$ have the same characteristic polynomials and hence they have the same eigenvalues with the same algebraic multiplicities. 
           \end{proof}
           
            \Answer 
            \begin{enumerate}
            \item {\bf False.} Let $A=\begin{bmatrix} 0 & 0 \\ 0 & 0\\ \end{bmatrix}$ and let $B=\begin{bmatrix} 0 & 1 \\ 0 & 0\\ \end{bmatrix}$. Then both $A$ and $B$ are upper triangular, and hence the eigenvalues are given by the entries on the main diagonal. Thus both $A$ and $B$ have eigenvalue 0 with algebraic multiplicity 2. However, for any invertible matrix $P$, we get $PAP^{-1}=A$. Thus $A$ and $B$ cannot be similar. 
            \item {\bf False.} Let $A=\begin{bmatrix} 0 & 1 \\ 0 & 0\\ \end{bmatrix}$ and let $B=\begin{bmatrix} 0 & 0 \\ 0 & 1\\ \end{bmatrix}$, then $A \sim B$, but they have different eigenvalues.
            \item {\bf False.} Again let $A=\begin{bmatrix} 0 & 1 \\ 0 & 0\\ \end{bmatrix}$ and let $B=\begin{bmatrix} 0 & 0 \\ 0 & 1\\ \end{bmatrix}$, then $A\sim B$. But they have different eigenvalues, so by the contrapositive of problem 9, we get $A$ and $B$ cannot be similar. 
            \end{enumerate}

\end{enumerate}  
\end{comment}

\end{document}


%%% Local ables: 
%%% mode: latex
%%% TeX-master: t
%%% End: 
